\documentclass[10pt]{article}
\usepackage[T1]{fontenc}
\usepackage[utf8]{inputenc}
\usepackage[french]{babel}
\usepackage{amsmath,amssymb,mathrsfs,mathpazo}
\usepackage{geometry}
\geometry{a4paper,margin=2cm}
\usepackage{graphicx,array,tabularx,tabularray,colortbl,multirow,multicol,booktabs,enumitem,pifont,ulem,fancybox,diagbox,cancel}
\usepackage{pstricks,pstricks-add,pst-plot,pst-tree,pst-node,pst-text,pst-func,pst-math,pst-bspline,pst-3dplot,pst-eucl,pst-solides3d}
\let\pscalculate\undefined
\usepackage{pst-bezier}
\usepackage{tikz,pgfplots,tkz-tab}
\usetikzlibrary{arrows,arrows.meta,patterns,calc,positioning,trees,decorations.markings,shapes.geometric,shapes.misc,shapes.arrows,intersections,decorations.pathmorphing,decorations.pathreplacing,backgrounds,shadows,babel}
\pgfplotsset{compat=1.18}
\usepgfplotslibrary{fillbetween}
\usepackage[np]{numprint}
\usepackage{listings,makecell,pas-tableur,twoopt}
\let\barre\undefined
\usepackage[pstricks]{tablvar}
\definecolor{darkgreen}{RGB}{0,100,0}
\providecommand{\red}{\color{red}}
\providecommand{\blue}{\color{blue}}
\providecommand{\green}{\color{green}}
\providecommand{\black}{\color{black}}
\providecommand{\white}{\color{white}}
\providecommand{\gray}{\color{gray}}
\providecommand{\cyan}{\color{cyan}}
\providecommand{\magenta}{\color{magenta}}
\providecommand{\yellow}{\color{yellow}}
\providecommand{\violet}{\color{violet}}
\providecommand{\orange}{\color{orange}}
\providecommand{\euro}{\texteuro}
\providecommand{\EUR}{\texteuro}
\providecommand{\ieme}{\textsuperscript{e}}
\providecommand{\degres}{\ensuremath{^\circ}}
\providecommand{\degree}{\ensuremath{^\circ}}
\providecommand{\vv}[1]{\overrightarrow{#1}}
\providecommand{\Coord}[1]{#1}
\providecommand{\decosix}{}
\providecommand{\starredbullet}{\ensuremath{\bullet}}
\pagestyle{empty}
\setlength{\parindent}{0pt}
\usepackage{adjustbox}\DeclareUnicodeCharacter{2236}{:}\DeclareUnicodeCharacter{2190}{\ensuremath{\leftarrow}}
\providecommand{\R}{\mathbb{R}}
\providecommand{\N}{\mathbb{N}}
\providecommand{\Z}{\mathbb{Z}}
\providecommand{\Q}{\mathbb{Q}}
\providecommand{\D}{\mathbb{D}}
\providecommand{\C}{\mathbb{C}}
\providecommand{\e}{{\mathrm{e}}}
\providecommand{\vect}[1]{\overrightarrow{#1}}
\providecommand{\vectt}[1]{\overrightarrow{\mathrm{#1}}}
\providecommand{\vvt}[1]{\overrightarrow{\mathrm{#1}}}
\providecommand{\np}[2][]{{#2~#1}}
\providecommand{\vv}[1]{{\overrightarrow{#1}}}
\renewcommand{\d}{{\mathrm{d}}}
\providecommand{\degres}{{^\circ}}
\providecommand{\degre}{{^\circ}}
\providecommand{\blue}{\color{blue}}
\providecommand{\red}{\color{red}}
\providecommand{\enskip}{\quad}
\providecommand{\llbracket}{\lbrack\!\lbrack}
\providecommand{\rrbracket}{\rbrack\!\rbrack}
\providecommand{\euro}{{\text{€}}}
\providecommand{\no}{{\text{n°}}}
\providecommand{\interval}[2]{{[#1;#2]}}
\providecommand{\intervFO}[2]{{[#1;#2[}}
\providecommand{\dotfill}{\ldots}

\providecommand{\ds}{\displaystyle}
\providecommand{\textding}[1]{#1}
\providecommand{\Oij}{\(\mathrm{O};\vec{\imath},\vec{\jmath}\)}
\providecommand{\Oijk}{\(\mathrm{O};\vec{\imath},\vec{\jmath},\vec{k}\)}
\providecommand{\Ouv}{\(\mathrm{O};\vec{u},\vec{v}\)}
\providecommand{\intervFF}[2]{\ensuremath{[#1;#2]}}
\renewcommand{\euro}{\ifmmode\text{\texteuro}\else\texteuro\fi}
\DeclareUnicodeCharacter{00B0}{\ifmmode^{\circ}\else\textdegree\fi}
\DeclareUnicodeCharacter{20AC}{\ifmmode\text{\texteuro}\else\texteuro\fi}
% Énoncés : Éducation nationale. Transcription et corrigés : APMEP, Denis Vergès et contributeurs. Découpage et classement : CoopMaths, Sébastien Lozano, Éric Elter et contributeurs. Adaptation de la présentation : Maths Explorer.
% CC BY-SA : https://coopmaths.fr/www/about/
% https://coopmaths.fr/alea/static/e3c/2020/tex/e3c_2020_00_sujet58_1.tex
\begin{document}
\noindent\textit{APMEP / CoopMaths. CC BY-SA. Présentation adaptée par Maths Explorer.}\par\medskip
\medskip

\emph{Ce QCM comprend $5$ questions indépendantes.}

\emph{Pour chacune d'elles, une seule des réponses proposées est exacte.}
 
\emph{Indiquer pour chaque question sur la copie la lettre correspondant à la réponse
choisie. Aucune justification n'est demandée.}

\emph{Chaque réponse correcte rapporte $1$ point. Une réponse incorrecte ou une absence de
réponse n'apporte, ni ne retire de point.}

\medskip

\begin{enumerate}
\item Pour tout réel $x$, $\e^{2x}+\e^{4x}$ est égal à

\begin{tabularx}{\linewidth}{*{4}{X}}
\textbf{a.~~} $\e^{6x}$ &\textbf{b.~~} $\e^{2x}(1+\e^{2})$&\textbf{c.~~}$\e^{3x}(\e^{x}+\e^{-x})$& \textbf{d.~~} $\e^{8x^2} $.
\end{tabularx}

\item Dans le plan muni d'un repère orthonormé \Oij, on considère les vecteurs $\vv{u}(-5~;~2)$ et \mbox {$\vv{v}(4~;~10)$} et la droite (d) d'équation : $5x + 2y + 3 = 0$.

\begin{tabularx}{\linewidth}{*{4}{X}}
\textbf{a.~~} $\vv{u}$ et $\vv{v}$  sont colinéaires &\textbf{b.~~} $\vv{u}$ est un vecteur normal à la droite (d) &\textbf{c.~~} $\vv{u}$ et $\vv{v}$  sont orthogonaux& \textbf{d.~~} $\vv{u}$ est un vecteur directeur de (d) .
\end{tabularx}

\item La dérivée $f'$ de la fonction $f$ définie sur $\R$ par $f(x)=(2x-1)\e^{-x}$  est : 

\begin{tabularx}{\linewidth}{*{4}{X}}
\textbf{a.~~}  $2x\e^{-x}$ &\textbf{b.~~}$-2x\e^{-x}$ &\textbf{c.~~}$(-2x+3)\e^{-x}$& \textbf{d.~~}$2\e^{-x}+(2x-1) \e^{-x}$  .
\end{tabularx}

\item Pour tout réel $x$, on a $\sin(\pi+x)=$

\begin{tabularx}{\linewidth}{*{4}{X}}
\textbf{a.~~} $-\sin (x)$ &\textbf{b.~~} $\cos (x)$&\textbf{c.~~}$\sin (x)$& \textbf{d.~~} $-\cos (x)$ .
\end{tabularx}
\smallskip

\begin{multicols}{2}
 \item Soit $f$ une fonction définie et dérivable sur $\R$ dont la courbe représentative est donnée ci-contre. La tangente à la courbe au point A est la droite T.


\psset{unit=0.95cm,labelFontSize=\scriptstyle,labelsep=0.1pt,showorigin=false}
\begin{pspicture}(-1.5,-2.8)(3.5,3.5)
\multido{\n=-1+1}{8}{\psline[linewidth=0.75pt,linecolor=lightgray](\n,-2.1)(\n,3.2)}
\multido{\n=-2+1}{6}{\psline[linewidth=0.75pt,linecolor=lightgray](-1.1,\n)(6.1,\n)}
\psaxes[linewidth=0.95pt,]{->}(0,0)(-1.2,-2.2)(6.1,3.3)
\psdots[dotstyle=+,dotscale =1.4,dotangle=45](0,3) 
\uput[dl](0,0){O} \uput[dr](0.2,3.2){A} \uput[u](1.7,-2){\blue $\mathcal{C}_f$}\uput[d](0.5,-1.1){\red T}

\psplot[linewidth=1pt,linecolor=red,plotpoints=3000]{0}{1}{x 5 neg mul 3 add}
\psplot[linewidth=1.25pt,linecolor=blue,plotpoints=2000]{-0.1}{6}{0.9 x sub 2.71828 x 0.5 mul exp div 3.333 mul}
\end{pspicture}
\end{multicols}

\begin{tabularx}{\linewidth}{*{4}{X}}
\textbf{a.~~}$f'(0)=3$  &\textbf{b.~~}$f'(0)=\frac{1}{5}$  &\textbf{c.~~}$f'(0)=5$& \textbf{d.~~}$f'(0)=-5$.
\end{tabularx}
 
\end{enumerate}

\vspace{0.75cm}

\end{document}
